% Generated draft: requires mathematical and source review.
\documentclass[11pt]{article}
\usepackage[margin=25mm]{geometry}
\usepackage{fontspec}
\setmainfont{Latin Modern Roman}
\usepackage{amsmath,amssymb,mathtools,unicode-math}
\setmathfont{Latin Modern Math}
\usepackage{enumitem,needspace,xurl,hyperref}
\setlength{\parindent}{0pt}
\setlength{\parskip}{5pt}
\setlength{\emergencystretch}{3em}
\newcommand{\N}{\mathbb{N}}
\newcommand{\Z}{\mathbb{Z}}
\newcommand{\Q}{\mathbb{Q}}
\newcommand{\R}{\mathbb{R}}
\newcommand{\C}{\mathbb{C}}
\newcommand{\F}{\mathbb{F}}
\newcommand{\A}{\mathbb{A}}
\newcommand{\PP}{\mathbb{P}}
\newcommand{\E}{\mathbb{E}}
\newcommand{\eps}{\varepsilon}
\newcommand{\dd}{\,\mathrm{d}}
\begin{document}
\section*{UnsolvedMath ID: 2306053}
\textbf{Source problem code:} catalogue entry
\subsection{Definitions and mathematical statement}
A domain $G\subset\C$ is convex in the direction of the imaginary axis (CIA) if, for every $x\in\R$, the intersection $G\cap\{x+iy:y\in\R\}$ is connected or empty. Let $\mathcal C$ consist of univalent functions
\[
f(z)=z+a_2z^2+\cdots\quad(z\in\mathbb D)
\]
whose image $G_1=f(\mathbb D)$ is CIA, and set $G_r=f(\{|z|<r\})$ for $0<r<1$.

D. A. Brannan's four questions are:
\begin{enumerate}
\item Determine the universal supremum
\[
r_* =\sup\{r\in(0,1):G_r\text{ is CIA for every }f\in\mathcal C\}.
\]
The source says a result of Goodman--Saff suggests $r_*=\sqrt2-1$; this value is a suggestion, not an asserted theorem. The endpoint $r=1$ is excluded from the supremum because $G_1$ is CIA by hypothesis.
\item Find reasonable sufficient conditions on $G_1$, equivalently on its normalized conformal map $f$, ensuring $G_r$ is CIA for every $r\in(0,1)$.
\item Slit $G_1$ along the real axis, and take the two components $G',G''$ adjacent to $0$, with $0$ on their boundaries. The source asks whether $G_1=G'\cup G''$ implies that every $G_r$ is CIA. This premise is inconsistent under the literal slit-domain interpretation: the components omit the real axis, in particular $0$, whereas $0=f(0)\in G_1$. Whether the intended condition uses the slit domain, closures, or restored slit intervals must be recovered; no correction is selected here.
\item If $f\in\mathcal C$, $0<r_0<1$, and $G_{r_0}$ is not CIA, must $G_r$ fail to be CIA for every $r_0<r<1$? This does not include the endpoint $1$.
\end{enumerate}
The source reports that $G_1$ being CIA need not make all $G_r$ CIA, nor provide a universal terminal radius range on which CIA is guaranteed. It attributes these negative observations to the authors transcribed as Hentgartner--Schobe and Goodman--Saff. The supplied summary reports openness as of the 2018 edition and no definitive resolution located. No proof or new sufficient condition is supplied here.
\subsection{Short English statement}
The questions concern when concentric subdisc images of a vertically convex univalent image remain vertically convex: a universal radius, sufficient conditions, a slit-domain criterion, and persistence of failure. The slit-domain criterion has a literal set-equality defect that must be checked in the source.
\subsection{Sources}
\begin{itemize}
\item Local UnsolvedMath registry, permanent ID 2306053, exact source code catalogue entry, ``Research Problems in Function Theory --- Problem 6.53''; status open, classification OPEN-TRIAGE.
\item Hayman, \emph{Research Problems in Function Theory} (2018), Problem 6.53, attributed to D. A. Brannan, \url{https://arxiv.org/abs/1809.07200}. The record identifies extraction from public arXiv source TeX, accessed July 29, 2026; the 2018 update was retained for status triage.
\item Hentgartner and Schobe (names as transcribed), and Goodman and Saff, credited for the failure of CIA under subdisc images. A Goodman--Saff result is also said to suggest $\sqrt2-1$ for the universal radius. No complete bibliographic references or precise theorem statements are supplied.
\item The supplied research summary reports openness as of the 2018 edition and no definitive resolution located, with no additional authoritative citation or precise search date.
\end{itemize}
\end{document}
